\section{Method}
\label{sec:method}

We study a single, precise question: does dialogue context rotate the low-dimensional
residual-stream subspaces that encode hallucination/epistemic state and safety/refusal,
beyond what is explained by trivial nuisance variables (prompt length, final-token
position, and surface topic)? This section defines the estimator, the gauge-invariant
drift metric, the projector-leak identity that operationalizes ``a static intervention is
mis-specified,'' the three confounds and their controls, and the pre-registered decision
rule. Throughout, a residual activation is $x \in \mathbb{R}^d$, read at the last token of
a prompt at a fixed layer.

\subsection{Subspace estimation by contrast}
\label{sec:subspace}

At dialogue turn $t$ we form a prompt by prepending the running conversation (the first
$t$ user turns, each followed by a fixed, neutral assistant reply that carries no
task-relevant content) and appending a held-out contrastive probe as the next user turn.
The conversation context is the \emph{only} thing that varies across turns: the probe bank
is fixed across all turns and all dialogues, so any movement we measure in the recovered
subspace is attributable to the prepended context rather than to a change of probes.

We recover concept subspaces by contrastive difference-of-means, the standard construction
behind linear probing and steering directions
\citep{marks2023geometry, arditi2024refusal, panickssery2024caa}. For the
\emph{hallucination/epistemic} subspace $H_t$ we use answerable real-entity questions
paired with surface-matched unanswerable fictional-entity questions, the
knowledge-awareness contrast of \citet{ferrando2025entity}; for the \emph{safety} subspace
$S_t$ we use harmful requests paired with surface-matched harmless requests, the refusal
contrast of \citet{arditi2024refusal}. Let $\{(x_i^{+}, x_i^{-})\}_{i=1}^{n}$ be the
last-token activations of the $n$ matched pairs at turn $t$, and let
$\delta_i = x_i^{+} - x_i^{-}$ be the per-pair difference vectors. Stacking the (centered)
$\delta_i$ into $D \in \mathbb{R}^{n \times d}$, we take $U_t \in \mathbb{R}^{d \times k}$
to be an orthonormal basis for the top-$k$ right singular vectors of $D$ (equivalently the
top-$k$ eigenvectors of $\sum_i \delta_i \delta_i^\top$). The default is the rank-1
difference-of-means direction, $k=1$, normalized to unit length;
Section~\ref{sec:power} explains why higher rank is not usable at our probe budgets. Pairing
removes the content the two sides share by construction, so $U_t$ isolates the concept axis
(epistemic status, or harmfulness) rather than the surface form of the probe.

\subsection{Gauge-invariant chordal drift}
\label{sec:drift}

A difference-of-means subspace is recovered only up to an arbitrary orthonormal basis: the
estimator returns \emph{a} basis $U_t$ for the plane $\mathrm{span}(U_t)$, and any
$U_t R$ with $R \in O(k)$ spans the same plane. A drift metric must therefore be a
function of the planes, not of the bases. Given two $k$-planes with orthonormal bases
$U, V \in \mathbb{R}^{d \times k}$, let $\theta_1, \dots, \theta_k$ be their principal
angles, defined through the singular values $\sigma_i = \cos\theta_i$ of $U^\top V$
\citep{edelman1998geometry}. We measure drift by the normalized chordal distance
\begin{equation}
\label{eq:chordal}
d(U, V) \;=\; \sqrt{\frac{1}{k} \sum_{i=1}^{k} \sin^2 \theta_i} \;\in\; [0, 1],
\end{equation}
which is $0$ when the planes coincide ($\theta_i = 0$ for all $i$) and $1$ when they are
orthogonal ($\theta_i = \pi/2$).

\paragraph{Derivation of the projector form and basis-invariance.} The quantity in
\eqref{eq:chordal} equals a Frobenius distance between the two orthogonal projectors
$P_U = U U^\top$ and $P_V = V V^\top$, which makes its basis-invariance manifest. Using
$\|A\|_F^2 = \mathrm{tr}(A^\top A)$ and that projectors are symmetric idempotents
($P_U^\top P_U = P_U$, $\mathrm{tr}\,P_U = k$),
\begin{equation}
\label{eq:projexpand}
\| P_U - P_V \|_F^2
= \mathrm{tr}(P_U) + \mathrm{tr}(P_V) - 2\,\mathrm{tr}(P_U P_V)
= 2k - 2\,\mathrm{tr}(U U^\top V V^\top).
\end{equation}
The cross term is $\mathrm{tr}(U U^\top V V^\top) = \| U^\top V \|_F^2 = \sum_i \sigma_i^2
= \sum_i \cos^2\theta_i$, the sum of squared principal-angle cosines. Substituting and
using $\sin^2\theta_i = 1 - \cos^2\theta_i$,
\begin{equation}
\label{eq:projid}
\| P_U - P_V \|_F^2 = 2k - 2 \sum_{i} \cos^2\theta_i = 2 \sum_{i} \sin^2\theta_i
\quad\Longrightarrow\quad
d(U, V) = \frac{1}{\sqrt{2k}}\, \| U U^\top - V V^\top \|_F .
\end{equation}
The right-hand side depends only on the projectors $P_U, P_V$, which are themselves
invariant to the choice of basis within each plane ($U R$ gives $U R R^\top U^\top =
U U^\top$ for $R \in O(k)$). Hence $d(U, V)$ is exactly the gauge-invariant principal-angle
quantity, and substituting an arbitrary estimator basis cannot change it. We also use the
cross-subspace overlap
\begin{equation}
\label{eq:omega}
\Omega(S_t, H_t) \;=\; \frac{1}{k}\, \big\| U_S^\top U_H \big\|_F^2 \;\in\; [0, 1],
\end{equation}
(with $U_S = U_t^{S}$, $U_H = U_t^{H}$) where $\Omega = 1$ means the safety and
hallucination planes coincide and $\Omega = 0$ means they are orthogonal; $\Omega$ is the
quantity a static orthogonalizer drives to zero once, and by the same projector argument it
is basis-invariant.

\subsection{The static-projector leak identity}
\label{sec:leak}

The operational claim behind conditional disentanglement is that a \emph{static}
intervention learned at one turn is mis-specified at later turns. We make this measurable
with an exact identity. Suppose at turn $1$ we fit a static projector that nulls the
hallucination subspace, $\Pi_1 = I - U_1 U_1^\top$ with $U_1 = U_1^{H}$ (orthonormal,
$k=1$ by default). At a later turn $t$ the concept has moved to $H_t = \mathrm{span}(U_t)$,
and the \emph{re-entanglement leak} is the fraction of a unit vector along $U_t$ that the
frozen projector fails to remove. The surviving component is
$\Pi_1 U_t = U_t - U_1 (U_1^\top U_t)$, with squared norm
\begin{equation}
\label{eq:leak}
\| \Pi_1 U_t \|_F^2 = \| U_t \|_F^2 - \| U_1^\top U_t \|_F^2
= k - \sum_i \cos^2 \theta_i = \sum_i \sin^2 \theta_i,
\end{equation}
so the root-mean-square surviving fraction is
$\sqrt{\tfrac{1}{k}\sum_i \sin^2\theta_i} = d(U_1, U_t) = d(H_1, H_t)$, exactly the chordal
drift of \eqref{eq:chordal}. In words, the static projector that nulls $H_1$ leaves a
residual on $H_t$ whose RMS magnitude is precisely the principal-angle drift between the
two subspaces, because the surviving fraction along each principal direction is the sine of
its principal angle. A conditional projector rebuilt at turn $t$ has leak $0$ by
construction. This identity lets us evaluate a static baseline's degradation directly from
the drift metric, and crucially lets us subject that degradation to the same
matched-null control as the drift itself (Section~\ref{sec:controls}).

\subsection{Three confounds and their controls}
\label{sec:controls}

``The subspace moved with context'' is easy to measure badly. Three confounds dominate a
naive measurement, and each, left uncontrolled, produces a false positive
(Section~\ref{sec:cascade}). Each has an explicit control.

\paragraph{(1) No-context onset.} Turn $0$ has no prepended context, so the cumulative
drift $d(U_0, U_t)$ is dominated by the generic transition from a bare probe to a
probe-in-context, a large displacement that is shared across all dialogues and carries no
content-conditional information. \emph{Control:} take the baseline at turn $1$ (the first
in-context turn) and measure $d(U_1, U_t)$, reporting the onset $d(U_0, U_1)$ separately.

\paragraph{(2) Token length and final-token position.} As a dialogue accumulates, the
prompt grows monotonically, and the last-token activation of a position-encoded model
shifts with absolute prompt length even at fixed semantic content. \emph{Control:} a
\textbf{scrambled-twin null}. For each user turn we word-shuffle the tokens, producing a
twin that preserves the \emph{exact} token multiset (hence identical length, lexical
content, and bag-of-words topic) while destroying coherent meaning. Genuine
content-conditionality requires the real subspace to move \emph{more} than its
length/topic-matched scrambled twin; if the two move equally, the movement is a length/topic
artifact rather than a content effect. The scrambled twin is the load-bearing control of
this study.

\paragraph{(3) Single-dialogue idiosyncrasy.} A between-dialogue spread computed from a
handful of dialogues can be carried by a single topically-divergent script.
\emph{Control:} a dialogue-level bootstrap confidence interval, resampling whole dialogues
(not just probes within a dialogue), so the reported uncertainty reflects generalization
across dialogues.

\paragraph{Why the same-context bootstrap floor is the wrong null.} A tempting null is to
draw two independent bootstrap resamples of the \emph{same} probe set at a \emph{fixed}
context and call their drift the noise floor, then declare an effect whenever the
between-dialogue spread exceeds that floor. This is the wrong comparison for a
between-context statistic, and the error is structural rather than a matter of sample size.
The between-dialogue spread is a function of variation \emph{across} contexts: it contains
all of the length, position, and topic variance that differs from one dialogue (or turn) to
another. The same-context floor, by construction, resamples probes within one fixed context
and therefore contains \emph{zero} between-context variance: it measures only
probe-estimation noise. Comparing a statistic that carries between-condition nuisance
variance against a floor that carries none makes ``spread $>$ floor'' near-tautological; it
will fire on length and topic alone, with no coherent content involved. The correct null is
one that \emph{matches} the nuisance variables (length, position, topic) and differs only in
the variable of interest (coherent meaning), which is exactly what the scrambled twin
provides. This single substitution, same-context floor $\rightarrow$ scrambled-twin null, is
what converts our headline from positive to null (Section~\ref{sec:cascade}).

\subsection{Dialogue-level bootstrap and the pre-registered decision rule}
\label{sec:rule}

For each layer and each condition (real, scrambled) we compute the median pairwise
between-dialogue chordal spread of the final-turn (full-context) subspaces, and bracket it
with a $95\%$ confidence interval obtained by resampling dialogues with replacement
(dialogue-level bootstrap). The decision rule is fixed in advance:
\begin{quote}
A layer is \emph{content-conditional} iff the real between-dialogue spread's CI lower bound
exceeds the scrambled-twin spread's CI upper bound. The model/space is judged
\textsc{content\_conditional} iff there is a contiguous band of $\geq 3$ such layers;
otherwise the verdict is \textsc{not\_content\_conditional}.
\end{quote}
The contiguity requirement guards against isolated layers that clear the bar by chance:
with per-layer false-positive risk near the nominal level, the expected number of
single-layer crossings over a network is non-negligible, but a run of three contiguous
crossings is not. We report the realized per-layer comparisons and the floor regardless of
the verdict, so a null is never confused with an under-powered run.

\subsection{Validation against analytic ground truth}
\label{sec:validation}

We validate the estimator and metric on synthetic activations with a known contrast
subspace that we rotate by a prescribed angle per turn. Three checks pass. (i)
\emph{Principal-angle sanity:} identical subspaces return $d = 0$ and orthogonal subspaces
return principal angles of $\pi/2$ (so $d = 1$), confirming the implementation of
\eqref{eq:chordal}. (ii) \emph{Recovery of an injected rotation:} a $15^\circ$-per-turn
injected rotation is recovered as noise-corrected drift $0.144$, against the analytic value
$0.149$. (iii) \emph{Correct negative:} a static (zero-rotation) subspace returns drift at
the bootstrap noise floor, i.e.\ a correctly \textsc{negative} verdict. The same code paths
are used for the real-model runs, so the protocol is calibrated to call both a true effect
and a true null.

\subsection{Statistical power and the rank/probe budget}
\label{sec:power}

Subspace estimation in $d \approx 10^3$ activation space is noisy, and the noise floor of
the chordal drift scales like $1/\sqrt{n_{\text{probes}}}$. On
\texttt{Qwen2.5-0.5B-Instruct} (raw activations, rank-1) the $95\%$ bootstrap floor falls
from $0.49$ at $16$ probes, to $0.31$ at $32$, $0.25$ at $64$, and $0.22$ at $96$ probes
(Table~\ref{tab:power}). Higher rank is not usable at these budgets: rank-2 and rank-3
subspaces estimated from a few dozen probes have a floor near $0.7$, meaning two bootstrap
estimates of the \emph{same} subspace come out nearly orthogonal, i.e.\ pure noise. We
therefore fix the default to rank-1 difference-of-means directions with $40$--$96$ probe
pairs per side, and we report the floor alongside every result so that a null is never
mistaken for an under-powered run.

\begin{table}[t]
\centering
\small
\begin{tabular}{@{}lccccc@{}}
\toprule
$n_{\text{probes}}$ & 16 & 32 & 64 & 96 & scaling \\
\midrule
rank-1 floor (95\%) & 0.49 & 0.31 & 0.25 & 0.22 & $\sim 1/\sqrt{n_{\text{probes}}}$ \\
\bottomrule
\end{tabular}
\caption{Bootstrap noise floor of the rank-1 chordal drift on
\texttt{Qwen2.5-0.5B-Instruct} (raw activations) as a function of the number of probe pairs.
The floor scales as $1/\sqrt{n_{\text{probes}}}$. Rank-2/3 from a few dozen probes sits at a
floor $\approx 0.7$ (pure noise), so we use rank-1 with $40$--$96$ probes.}
\label{tab:power}
\end{table}

\section{Experimental Setup}
\label{sec:setup}

\paragraph{Models and representations.} We evaluate three settings, chosen to separate the
null from any single model scale or activation basis. (i) \texttt{Qwen2.5-0.5B-Instruct}
\citep{qwen25}, raw residual activations, $24$ layers (cached $896$-dimensional last-token
states). (ii) \texttt{Qwen2.5-1.5B-Instruct} \citep{qwen25}, raw residual activations, $28$
layers. (iii) \texttt{GPT-2} small in \emph{SAE feature space}, using the publicly released
\texttt{gpt2-small-res-jb} residual-stream sparse autoencoders at all $12$ layers
\citep{cunningham2023sparse, bricken2023monosemanticity}: each contrastive probe's
last-token activation is encoded through the frozen SAE and the subspace is recovered in
feature coordinates. The third setting is the one that matters for the methods we test:
conditional disentanglement operates on SAE features, not raw activations
\citep{mahmoud2025unintended}, so a null in raw space alone would be inconclusive. We read
the last-token activation \emph{per layer}, indexing the last token before stacking hidden
states, and recover $U_t$ at each layer independently.

\paragraph{Contrastive probe construction.} The hallucination/epistemic probe bank pairs
answerable real-entity questions with surface-matched unanswerable fictional-entity
questions \citep{ferrando2025entity}; the matching holds the question template, syntax, and
length fixed across the pair, so the difference vector isolates epistemic status rather than
surface form. The safety probe bank pairs harmful requests with surface-matched harmless
requests \citep{arditi2024refusal}, matched the same way. We use rank-1 difference-of-means
directions and $40$--$96$ probe pairs per side (Section~\ref{sec:power}), reading the
last-token residual at every layer. The probe banks are fixed across all turns and
dialogues, so the only thing that varies within a dialogue is the prepended context.

\paragraph{Dialogues and the scrambled twin.} We use $12$ topically diverse multi-turn
dialogues (five user turns each, each followed by a fixed neutral assistant reply). For
each real dialogue we construct its \textbf{scrambled twin} by word-shuffling each user
turn, which holds the token multiset (length, lexicon, and bag-of-words topic) exactly fixed
while destroying coherent meaning (Section~\ref{sec:controls}). Twelve dialogues is the
operative sample for the dialogue-level bootstrap; the rigor cascade
(Section~\ref{sec:cascade}) shows that the smaller $n{=}3$ placebo is under-supported and can
manufacture a spurious mid-layer effect, which is why the controlled runs use $n{=}12$.

\paragraph{Headline statistic.} For each setting we take, at every layer, each dialogue's
final-turn (full-context) hallucination subspace, compute the median pairwise
between-dialogue chordal spread, bracket it with a dialogue-level $95\%$ bootstrap CI, and
compare it against the scrambled-twin spread CI under the decision rule of
Section~\ref{sec:rule}. For the static-projector demonstration we instead report the leak
$d(H_1, H_t)$ of \eqref{eq:leak} against its scrambled twin, under the same rule. The
safety subspace $S_t$ is recovered identically; on the models we evaluate it is poorly
recoverable below the mid layers (its same-subspace floor reaches $\approx 1.0$), so
$\Omega(S_t, H_t)$ is meaningful only at late layers, a scope limitation we state with each
safety-side result. All numbers reported in this paper are from real runs on the released
harness; every figure and number is regenerated by the released scripts, with raw
activations and SAE features differing only by a backend swap.